\documentclass[journal,10pt]{IEEEtran}
\setlength{\textwidth}{182mm}
\setlength{\columnsep}{4.2mm}
\usepackage{amsmath,amssymb,newtxmath}
\usepackage[no-math]{fontspec}
\setmainfont{Times New Roman}
\usepackage{xeCJK}
\setCJKmainfont[AutoFakeBold=2]{SimSun}
\usepackage{graphicx,booktabs,array,hyperref,placeins}
\hypersetup{hidelinks,unicode=true}
\raggedbottom
\setlength{\emergencystretch}{2em}
\renewcommand{\arraystretch}{1.15}
\usepackage{seqsplit}
\begin{document}
\title{Litz 线分层建模与数值验证：完整研究报告 · 09 Ansys 算例与接口验证}
\author{缪雨轩\thanks{淮南师范学院}}

\maketitle
\noindent 合并稿对应章节 / Chapter in merged manuscript: \href{report_zh-CN_ieee.pdf#chapter-09}{09}\par
\FloatBarrier
\section*{09 Ansys 算例与接口验证}
\hypertarget{chapter-09}{}
\pdfbookmark[0]{09 Ansys 算例与接口验证}{bookmark-chapter-09}
七股二维历史有效端口短测使用 704 单元，包含两个电压基解和实际 ±20 A 场解；公共端电压、逐股复电流与独立功率支持该单档算例。它不是三档网格收敛。新固定网格与精度尝试未形成要求的完整验收；原生三维 J 表达式探针也未得到可接受物理解，因此不为其生成场结果图或填造数值。


归档边界采用已核查 v252 Student 的二维 2000 三角形、三维 64000 体单元/8000 面单元及四核限制。接口错误、网格或许可限制和数值误差有不同含义，应分别记录。实际端接、曲线 STL 和足铜大束继续需要相应许可、网格能力及独立验证；有限短测不能用于推荐真实换位性能。


\begin{figure}[!htbp]
\centering
\setcounter{figure}{17}
\refstepcounter{figure}\label{fig:ansys_seven_current}
\hypertarget{fig-ansys_seven_current}{}
\includegraphics[width=88.9mm]{../../../panels/zh-CN/ansys_seven_current.pdf}
\par\vspace{4pt}{\small\centering 图 18. Ansys 七股二维复电流。\par}
\vspace{3pt}{\small\raggedright 704 单元有效端口短测给出逐股幅值与相位。 电流采用实际 ±20 A 场解的端口记录。幅值与相位对应同一七股编号。\par}
\end{figure}
\begin{figure}[!htbp]
\centering
\setcounter{figure}{18}
\refstepcounter{figure}\label{fig:ansys_limits}
\hypertarget{fig-ansys_limits}{}
\includegraphics[width=88.9mm]{../../../panels/zh-CN/ansys_limits.pdf}
\par\vspace{4pt}{\small\centering 图 19. Ansys Student 能力与接口诊断。\par}
\vspace{3pt}{\small\raggedright Ansys Student v252 的二维三角形、三维体单元与面单元上限分别为 2000、64000 和 8000，计算使用最多 4 核。并列分组表示各项许可限制。电流密度表达式 J 的接口探针与许可规模分别记录，实际求解状态见三维诊断。\par}
\end{figure}
\end{document}
